% !PRINTED 250
\weberpage{250}
%<<CONT
zwei imaginären Wurzeln dar, die Punkte von $[2, 4]$ Gleichungen mit zwei gleichen und zwei davon verschiedenen reellen Wurzeln.

Auf der Doppelparabel finden zweimal zwei gleiche Wurzeln statt, und zwar auf dem Theil, in dem $[0]$ an $[4]$ grenzt, reelle, in dem isolirten Theil imaginäre.

% Source figure: Weber_Algebra_Bd1_IA.pdf, PDF page 276, printed p.250, Fig.5; exact caption-inclusive crop.
% Former placeholder description retained: perspective line drawing of the discriminant surface, showing two cuspidal-edge branches meeting at a shaded curved vault-like shell; dotted lines indicate hidden edges and the coordinate framework of the model.
\weberfigureimage{b1_i0275_fig5_source_20260831.png}{0.56}

Die Punkte der Rückkehrkanten repräsentiren Gleichungen mit drei gleichen Wurzeln und der Coordinaten-Anfangspunkt die Gleichung mit vier gleichen Wurzeln.

Auf die Discriminantenfläche und ihre Bedeutung für die Discussion der biquadratischen Gleichung hat zuerst \gesperrt{Kronecker} hingewiesen. Ein Modell der Fläche ist von \gesperrt{Kerschensteiner} construirt.\footnote{\gesperrt{Kronecker}, Monatsbericht der Berliner Akademie vom 14. Februar 1878. Eine Beschreibung des Modells findet sich in dem von Dyck herausgegebenen Katalog mathematischer Modelle, München 1892, dem die oben stehende Fig. 5 entnommen ist.}

Der analytische Ausdruck für die Bedingungen der Realität der Wurzeln lässt sich in eine Gestalt bringen, in der nur die Covarianten der biquadratischen Form vorkommen.\footnote{\gesperrt{Clebsch}, ,,Theorie der binären Formen`` \S.\,47. \gesperrt{Faà di Bruno}, ,,Theorie der binären Formen``, deutsch von Walter (Leipzig 1881), \S.\,20, woselbst sich ein wesentlicher Zusatz von \gesperrt{Nöther} findet.}

Den Ausgangspunkt dazu bilden die Ausdrücke für die Functionen $\psi_1$, $\psi_2$, $\psi_3$, die \S.\,66, (6) gegeben sind. Wenn wir in
%CONT>>
